% Part: normal-modal-logic
% Chapter: tableaux
% Section: simple-S5

\documentclass[../../../include/open-logic-section]{subfiles}

\begin{document}

\olfileid[hi]{nml}{tab}{s5}

\olsection{\Log{S5} के लिए सरल \usetoken{P}{tableau}}

\Log{S5} सार्विक प्रतिरूपों के वर्ग के सापेक्ष यथार्थ और पूर्ण है, अर्थात्
ऐसे प्रतिरूपों के वर्ग के सापेक्ष जिनमें प्रत्येक जगत प्रत्येक जगत से
अभिगम्य है। सार्विक प्रतिरूपों में अभिगम्यता संबंध से अंतर नहीं पड़ता:
``ऐसा जगत~$w$ है जहाँ $\mSat{M}{!A}[w]$'' तभी सत्य है जब ऐसा $w$ है जो
$u$ से अभिगम्य है। अतः \Log{S5} में हम प्रतिरूप को केवल जगतों के समुच्चय
और मूल्यांकन~$V$ के रूप में परिभाषित कर सकते हैं। इससे संकेत मिलता है कि
हम !!{tableau} नियमों को भी सरल कर सकते हैं। सामान्य स्थिति में हम धन
पूर्णांकों के अनुक्रमों को उपसर्ग लेते हैं, ताकि हिसाब रख सकें कि कौन-से
उपसर्ग ऐसे जगतों को नाम देते हैं जो दूसरों से अभिगम्य हैं: $\sigma.n$,
$\sigma$ से अभिगम्य जगत को नाम देता है। किंतु \Log{S5} में कोई भी जगत
किसी भी जगत से अभिगम्य है, अतः ऐसा हिसाब रखने की आवश्यकता नहीं। इसके
स्थान पर हम धन पूर्णांकों को उपसर्ग बना सकते हैं। सरल किए गए नियम
\olref{tab:rules-S5} में दिए गए हैं।

\begin{table}
  \begin{center}
    \def\arraystretch{3}\def\fCenter{}
    \begin{tabular}{|c|c|}
    \hline
    \iftag{prvBox}{
      \AxiomC{\sFmla{\True}{\Box !A}[n]}
      \RightLabel{\TRule{\True}{\Box}}
      \UnaryInfC{\sFmla{\True}{!A}[m]}
      \DisplayProof
      &
      \AxiomC{\sFmla{\False}{\Box !A}[n]}
      \RightLabel{\TRule{\False}{\Box}}
      \UnaryInfC{\sFmla{\False}{!A}[m]}
      \DisplayProof\\
      $m$ प्रयुक्त है & $m$ नया है
      \\[1ex]
      \hline}{}
    \iftag{prvDiamond}{
      \AxiomC{\sFmla{\True}{\Diamond !A}[n]}
      \RightLabel{\TRule{\True}{\Diamond}}
      \UnaryInfC{\sFmla{\True}{!A}[m]}
      \DisplayProof
      &
      \AxiomC{\sFmla{\False}{\Diamond !A}[n]}
      \RightLabel{\TRule{\False}{\Diamond}}
      \UnaryInfC{\sFmla{\False}{!A}[m]}
      \DisplayProof\\
      $m$ नया है & $m$ प्रयुक्त है
      \\[1ex]
      \hline}{}
    \end{tabular}
  \end{center}
  \caption{\Log{S5} के सरल किए गए नियम।}
  \ollabel{tab:rules-S5}
\end{table}

\begin{ex}
  हम ऐसा सरल किया गया संवृत टैब्लो देते हैं जो $\Log{S5} \Proves
  \Ax{5}$, अर्थात् $\Diamond!A \lif \Box\Diamond!A$, दिखाता है।
  \begin{oltableau}
    [\pFmla{\False}{\Diamond\formula{A} \lif \Box\Diamond \formula{A}}{1},
      just = \TAss
      [\pFmla{\True}{\Diamond \formula{A}}{1}, just = {\TRule{\False}{\lif}[1]}
        [\pFmla{\False}{\Box\Diamond \formula{A}}{1},
          just = {\TRule{\False}{\lif}[1]}
          [\pFmla{\False}{\Diamond \formula{A}}{2},
            just = {\TRule{\False}{\Box}[3]}
            [\pFmla{\True}{\formula{A}}{3},
              just = {\TRule{\True}{\Diamond}[2]}
                [\pFmla{\False}{\formula{A}}{3},
                  just = {\TRule{\False}{\Diamond}[4]}, close]
            ]
          ]
        ]
      ]
    ]
  \end{oltableau}
\end{ex}

\end{document}
